% !TeX program = lualatex
% ================================================================
% Polish Parallel Text Version of AlgebraicTerminology
%
% Generated by the server using the following settings:
%
%   language            Polish (pol)
%   text direction      left to right
%   font                Noto Serif
%   fallback font       Noto Serif
%   built from          latex/worksheets/algebra/algebraicTerminology.tex
%   vocabulary key      latex/worksheets/algebra/algebraicTerminology/L2/pol/algebraicTerminology_polishKey.tex
%   tier 3 vocabulary   green   RGB 0 128 0
%   English text        black   RGB 0 0 0
%   Polish text         purple  RGB 112 48 160
%   layout macros       version 3
%
% Please don't edit any information in this comment block - the server
% compares it against the current settings and rebuilds this file when
% they differ, so changes here would be overwritten. However, feel free
% to make updates to the LaTeX below if you want to edit it.
% ================================================================

\documentclass[14pt,a4paper]{extarticle}
% ================================================================
% Copied from the English sheet this was translated from, so that
% anything it sets up for itself still works here
% ================================================================
\usepackage[a4paper,top=12mm,bottom=12mm,left=14mm,right=14mm]{geometry}
\usepackage[T1]{fontenc}
\usepackage{amsmath}
\usepackage{amssymb}
% A sans serif font with a single-storey a, as children write it. It is
% used for the mathematics as well as the text.
\usepackage{tgadventor}
\renewcommand*{\familydefault}{\sfdefault}
\usepackage[italic]{mathastext}
\usepackage{tikz}
\usepackage[most]{tcolorbox}
\usetikzlibrary{arrows.meta,decorations.pathreplacing}

% ================================================================
% Algebraic terminology: a three-page sheet for a Year 7 class meeting
% algebra for the first time.
%
% Every page has the same two kinds of question. The first labels the
% parts of an expression, by writing their names in boxes. The second
% fills in the gaps in sentences about an expression, then goes the
% other way, from a description to the expression, and ends with an
% open task. The expressions have from 2 to 5 terms.
%
% Page 1 has the definitions and a labelled example of 5x+3, and there
% are no negative numbers. Some of the names are already in the boxes,
% and some of the gaps are filled in, so that pupils can make a start.
% Page 2 has negative constants and coefficients, and a coefficient of
% 1 or -1 that is not written. Page 3 adds decimals, squared terms, a
% constant that is not at the end, and an expression with no constant.
%
% The sheet is A4 in 14pt type, so that it can also be printed at A5
% (the type is then about 10pt). The boxes and gaps are sized to be
% written in at A5. The answers are in a comment beside each question.
% ================================================================

\pagestyle{empty}
\setlength{\parindent}{0pt}
\setlength{\parskip}{0pt}
% Ragged right, so that the spaces between words are all the same.
\AtBeginDocument{\raggedright}

% The idea the sheet is built on, boxed and only faintly tinted, so the
% sheet is cheap to print.
\newtcolorbox{keyidea}{colback=orange!3,colframe=orange!35,boxrule=0.6pt,
  left=6pt,right=6pt,top=4pt,bottom=4pt}

% An instruction for the questions after it.
\newcommand{\instruction}[1]{\par{\large\bfseries #1}\par\vspace{6pt}}

% A question number, and gaps to write in: \lgap for a letter, \gap for
% a number or a term, \tgap for a term with a decimal in it, and \wgap
% for a word or an expression. The strut keeps the lines far enough
% apart to write on at A5.
\newcommand{\qn}[1]{\textbf{Q#1.}\ }
\newcommand{\gapstrut}{\rule[-1.5mm]{0pt}{7.5mm}}
\newcommand{\lgap}{\gapstrut\underline{\hspace{1cm}}}
\newcommand{\gap}{\gapstrut\underline{\hspace{1.4cm}}}
\newcommand{\tgap}{\gapstrut\underline{\hspace{1.5cm}}}
\newcommand{\wgap}{\gapstrut\underline{\hspace{3.4cm}}}

% A gap that is filled in already, so that pupils can make a start:
% \given{$4$} for a number or a term, and \lgiven{$x$} for a letter.
\newcommand{\given}[1]{\gapstrut\underline{\makebox[1.4cm]{#1}}}
\newcommand{\lgiven}[1]{\gapstrut\underline{\makebox[1cm]{#1}}}

% An answer line at the right-hand end of a question. It goes on the
% last line of the question if there is room, and on a line of its own
% if not. (This is the \signed macro from The TeXbook, with \hfill in
% place of \hfil so that it still reaches the right margin when the
% text is ragged right.)
\newcommand{\answerline}[1]{{\unskip\nobreak\hfill\penalty50
  \hskip2em\hbox{}\nobreak\hfill\gapstrut\underline{\hspace{#1}}%
  \parfillskip=0pt \finalhyphendemerits=0 \par}}

% A labelling question: the number on the left, the picture centred.
\newcommand{\labelq}[2]{%
  \par\noindent\begin{minipage}[c]{1.3cm}\qn{#1}\end{minipage}\hfill
  #2\hfill\hspace*{1.3cm}\par}

% A fill-in-the-gaps question: the expression large, then the sentences.
\newenvironment{gapq}[3]
  {\begin{minipage}[t]{#1\linewidth}\raggedright
   \qn{#2}\quad{\Large $#3$}\par}
  {\end{minipage}}

% ================================================================
% Labelled expressions
% ================================================================

\tikzset{
  % The height and depth are in points, because an em here would be
  % measured in the body font and come out shorter than the letters.
  part/.style={inner xsep=0pt, inner ysep=2pt, font=\LARGE,
               text height=15pt, text depth=5pt},
  namebox/.style={draw, minimum width=3.8cm, minimum height=0.95cm,
                  inner sep=2pt},
  arr/.style={-{Stealth[length=7pt]}, thick, shorten >=3pt},
  brc/.style={decorate, decoration={brace, amplitude=6pt, raise=3pt}, thick},
  brcd/.style={decorate, decoration={brace, amplitude=6pt, raise=3pt,
               mirror}, thick},
}

% \anatomy[<box width>][<arrow length>]{<coefficient>}{<variable>}{<constant term>}
%         {<coefficient name>}{<variable name>}{<constant name>}{<expression name>}
% An expression of the form 5x+3 with an arrow to each part and a
% brace under the whole expression, each ending in a box. The box
% holds the name of the part, or is left empty for the name to be
% written in. The constant is written with its sign, because the sign
% in front of a term belongs to the term.
\NewDocumentCommand{\anatomy}{O{3.8cm}O{0.6cm}mmmmmmm}{%
  \begin{tikzpicture}[baseline=(v.base),
      namebox/.style={draw, minimum width=#1, minimum height=0.95cm,
                      inner sep=2pt}]
    \node[part] (c) {$#3$};
    \node[part, anchor=base west] (v) at (c.base east) {$#4$};
    \node[part, anchor=base west] (k) at ([xshift=9pt]v.base east) {$#5$};
    \node[namebox, anchor=east] (cl) at ([xshift=-#2]c.west) {#6};
    \draw[arr] (cl.east) -- (c.west);
    \node[namebox, anchor=south] (vl) at ([yshift=0.45cm]v.north) {#7};
    \draw[arr] (vl.south) -- (v.north);
    \node[namebox, anchor=west] (kl) at ([xshift=#2]k.east) {#8};
    \draw[arr] (kl.west) -- (k.east);
    \draw[brcd] (c.south west) -- (k.south east) coordinate[midway] (m);
    \node[namebox, anchor=north] at ([yshift=-0.4cm]m) {#9};
  \end{tikzpicture}%
}

% The other labelled expressions are built from the same pieces, so
% that a box can point at any part.
%
% \labelled{<first part>}{<parts and boxes>} draws the picture. Each
% part of the expression is a node with a name. \pfirst{c}{4} starts
% the expression with the part named c, \pnext{a}{c}{a} puts the part
% named a straight after c in the same term, and \pterm{t}{a}{+9b}
% starts a new term after a small gap.
%
% A box points at a part from the left of the first part (\bleft),
% the right of the last part (\bright), or above or below any part
% (\babove, \bbelow). \brabove and \brbelow put a brace over or under
% the parts from one named part to another, with a box beyond it.
\newcommand{\labelled}[2]{%
  \begin{tikzpicture}[baseline=(#1.base)]#2\end{tikzpicture}}
\newcommand{\pfirst}[2]{\node[part] (#1) {$#2$};}
\newcommand{\pnext}[3]{%
  \node[part, anchor=base west] (#1) at (#2.base east) {$#3$};}
\newcommand{\pterm}[3]{%
  \node[part, anchor=base west] (#1) at ([xshift=9pt]#2.base east) {$#3$};}
\newcommand{\bleft}[2]{%
  \node[namebox, anchor=east] (#1L) at ([xshift=-0.6cm]#1.west) {#2};
  \draw[arr] (#1L.east) -- (#1.west);}
\newcommand{\bright}[2]{%
  \node[namebox, anchor=west] (#1R) at ([xshift=0.6cm]#1.east) {#2};
  \draw[arr] (#1R.west) -- (#1.east);}
\newcommand{\babove}[2]{%
  \node[namebox, anchor=south] (#1A) at ([yshift=0.45cm]#1.north) {#2};
  \draw[arr] (#1A.south) -- (#1.north);}
\newcommand{\bbelow}[2]{%
  \node[namebox, anchor=north] (#1B) at ([yshift=-0.45cm]#1.south) {#2};
  \draw[arr] (#1B.north) -- (#1.south);}
\newcommand{\brabove}[3]{%
  \draw[brc] (#1.north west) -- (#2.north east) coordinate[midway] (bm);
  \node[namebox, anchor=south] at ([yshift=0.5cm]bm) {#3};}
\newcommand{\brbelow}[3]{%
  \draw[brcd] (#1.south west) -- (#2.south east) coordinate[midway] (bm);
  \node[namebox, anchor=north] at ([yshift=-0.4cm]bm) {#3};}

% Characters this language's font has not got are borrowed from another,
% rather than being silently dropped - see docs/translated-sheet-latex.md
\directlua{%
  if luaotfload and luaotfload.add_fallback then
    luaotfload.add_fallback("eallatinfallback",
      { "Noto Serif:mode=harf;" })
  else
    tex.error("this luaotfload is too old to register a fallback font")
  end
}
\usepackage[english]{babel}
\babelprovide[import]{polish}
\babelfont[polish]{rm}[Renderer=Harfbuzz, BoldFont={Noto Serif}, ItalicFont={Noto Serif}, BoldItalicFont={Noto Serif}, RawFeature={fallback=eallatinfallback}]{Noto Serif}
\babelfont[polish]{sf}[Renderer=Harfbuzz, BoldFont={Noto Serif}, ItalicFont={Noto Serif}, BoldItalicFont={Noto Serif}, RawFeature={fallback=eallatinfallback}]{Noto Serif}
\newcommand{\ealtext}[1]{\foreignlanguage{polish}{#1}}
\newcommand{\ealtextblock}[1]{{\raggedright\foreignlanguage{polish}{#1}\par}}
% ================================================================
% EAL layout helpers
% ================================================================
\usepackage{xcolor}

\definecolor{ealkeycolour}{RGB}{0,128,0}
\definecolor{ealenglish}{RGB}{0,0,0}
\definecolor{ealtwocolour}{RGB}{112,48,160}

\newsavebox{\ealboxen}
\newsavebox{\ealboxtr}
\newlength{\ealglosswd}
\newlength{\ealglossraise}
\setlength{\ealglossraise}{1.15em}

% a tier 3 word, in English and in translation alike
\newcommand{\ealkey}[1]{\textcolor{ealkeycolour}{#1}}
\newcommand{\ealkeytr}[1]{\textcolor{ealkeycolour}{\ealtext{#1}}}

% One translated word above one English word. Built from kernel
% primitives, never a tabular, which would break wherever the sheet
% loads array, booktabs or tabularx. The box is as wide as the wider of
% the two, so a long translation pushes its neighbours aside instead of
% overlapping them, and it claims its full height so a table row or a
% TikZ node makes room for it.
\newcommand{\ealgloss}[2]{%
  \sbox{\ealboxen}{\ealkey{#1}}%
  \sbox{\ealboxtr}{\small\textcolor{ealtwocolour}{\ealtext{#2}}}%
  \setlength{\ealglosswd}{\wd\ealboxen}%
  \ifdim\wd\ealboxtr>\ealglosswd\setlength{\ealglosswd}{\wd\ealboxtr}\fi
  \leavevmode
  \makebox[\ealglosswd][c]{%
    \makebox[0pt][c]{\usebox{\ealboxen}}%
    \makebox[0pt][c]{\raisebox{\ealglossraise}{\usebox{\ealboxtr}}}%
  }%
}

% opens up line spacing around a block of glosses, so the raised
% translations sit clear of the line above
\newenvironment{ealglossed}{\par\linespread{2.1}\selectfont\ignorespaces}{\par}

% a whole translated sentence above its English counterpart. block level,
% so it does not belong inside a tabular cell
\newcommand{\ealpara}[2]{%
  \par\addvspace{0.5\baselineskip}%
  {\color{ealtwocolour}\ealtextblock{#1}}%
  \penalty200
  {\color{ealenglish}#2\par}%
  \addvspace{0.5\baselineskip}%
}

\begin{document}

\begin{center}
  {\LARGE\bfseries \ealpara{Terminologia \ealkeytr{algebraiczna}}{\ealkey{Algebraic} Terminology}}
\end{center}

\begin{keyidea}
\begin{minipage}[c]{0.57\linewidth}
\raggedright
\ealpara{\ealkeytr{Zmienna} to litera, która może oznaczać dowolną liczbę.}{A \textbf{\ealkey{variable}} is a letter that could be any number.}

\vspace{3pt}
\ealpara{\ealkeytr{Wyraz} to liczba, \ealkeytr{zmienna} lub liczby i \ealkeytr{zmienne} pomnożone przez siebie. \ealkeytr{Znak} przed \ealkeytr{wyrazem} należy do tego \ealkeytr{wyrazu}.}{A \textbf{\ealkey{term}} is a number, a \ealkey{variable}, or numbers and \ealkey{variables} multiplied together. The \ealkey{sign} in front of a \ealkey{term} belongs to the \ealkey{term}.}

\vspace{3pt}
\ealpara{\ealkeytr{Współczynnik} \ealkeytr{wyrazu} to liczba, która mnoży \ealkeytr{zmienne} w nim.}{The \textbf{\ealkey{coefficient}} of a \ealkey{term} is the number that multiplies the \ealkey{variables} in it.}

\vspace{3pt}
\ealpara{\ealkeytr{Stała} to \ealkeytr{wyraz} bez \ealkeytr{zmiennej} w nim.}{A \textbf{\ealkey{constant}} is a \ealkey{term} with no \ealkey{variable} in it.}

\vspace{3pt}
\ealpara{\ealkeytr{Wyrazy} dodane lub odjęte, bez \ealkeytr{znaku} równości, tworzą \ealkeytr{wyrażenie}.}{\ealkey{Terms} added or subtracted, with no equals \ealkey{sign}, make an \textbf{\ealkey{expression}}.}
\end{minipage}\hfill
\begin{minipage}[c]{0.41\linewidth}
\centering
\tikzset{part/.append style={font=\Large}}%
\anatomy[1.9cm][0.3cm]{5}{x}{+3}{\small \ealkey{coefficient}}{\small \ealkey{variable}}{\small \ealkey{constant}}{\small \ealkey{expression}}
\end{minipage}
\end{keyidea}

\vspace{8pt}
\ealpara{Wpisz nazwę każdej części w jej pole.}{\instruction{Write the name of each part in its box.}}

\labelq{1}{\anatomy{4}{y}{+9}{}{\ealkey{variable}}{}{\ealkey{expression}}}

\vspace{14pt}
\labelq{2}{\labelled{c}{%
  \pfirst{c}{2}\pnext{a}{c}{a}\pterm{t}{a}{+6b}\pterm{k}{t}{+5}%
  \bleft{c}{}\bbelow{a}{}\brabove{t}{t}{\ealkey{term}}\bright{k}{}}}

\vspace{12pt}
\ealpara{Wypełnij luki.}{\instruction{Fill in the gaps.}}

\begin{gapq}{0.5}{3}{5xy+3z+7}
\ealpara{To \ealkeytr{wyrażenie} ma \gap\ \ealkeytr{wyrazy}.}{This \ealkey{expression} has \gap\ \ealkey{terms}.}

\ealpara{Jego \ealkeytr{zmienne} to \lgiven{$x$}, \lgiven{$y$} i \lgap.}{Its \ealkey{variables} are \lgiven{$x$}, \lgiven{$y$} and \lgap.}

\ealpara{\ealkeytr{Współczynnik} przy $xy$ wynosi \gap.}{The \ealkey{coefficient} of $xy$ is \gap.}

\ealpara{\ealkeytr{Stała} wynosi \gap.}{The \ealkey{constant} is \gap.}
\end{gapq}\hfill
\begin{gapq}{0.47}{4}{2m+9n+3p+4}
\ealpara{To \ealkeytr{wyrażenie} ma \given{$4$} \ealkeytr{wyrazy}.}{This \ealkey{expression} has \given{$4$} \ealkey{terms}.}

\ealpara{\ealkeytr{Współczynnik} przy $n$ wynosi \gap.}{The \ealkey{coefficient} of $n$ is \gap.}

\ealpara{\ealkeytr{Współczynnik} przy $p$ wynosi \gap.}{The \ealkey{coefficient} of $p$ is \gap.}

\ealpara{$4$ nazywa się \wgap.}{The $4$ is called the \wgap.}
\end{gapq}

\vspace{12pt}
\ealpara{\ealkeytr{Wyrażenie} ma \ealkeytr{zmienną} $p$. \ealkeytr{Współczynnik} przy $p$ wynosi $2$, a \ealkeytr{stała} wynosi $7$. Zapisz \ealkeytr{wyrażenie}.}{\qn{5} An \ealkey{expression} has the \ealkey{variable} $p$. The \ealkey{coefficient} of $p$ is $2$ and the \ealkey{constant} is $7$. Write down the \ealkey{expression}.}
\answerline{3.4cm}

\vspace{12pt}
\ealpara{Zapisz \ealkeytr{wyrażenie}, które ma $3$ \ealkeytr{wyrazy}, $2$ \ealkeytr{zmienne} i \ealkeytr{stałą} równą $8$.}{\qn{6} Write an \ealkey{expression} that has $3$ \ealkey{terms}, $2$ \ealkey{variables} and a \ealkey{constant} of $8$.}
\answerline{5.5cm}

\newpage

\ealpara{Wpisz nazwę każdej części w jej pole.}{\instruction{Write the name of each part in its box.}}

\labelq{7}{\anatomy{7}{y}{-2}{}{}{}{}}

\vspace{14pt}
\labelq{8}{\labelled{c}{%
  \pfirst{c}{4}\pnext{a}{c}{a}\pterm{t}{a}{+9b}\pterm{k}{t}{-6}%
  \bleft{c}{}\bbelow{a}{}\brabove{t}{t}{}\bright{k}{}}}

\vspace{14pt}
\labelq{9}{\labelled{c}{%
  \pfirst{c}{3}\pnext{p}{c}{p}\pnext{q}{p}{q}\pterm{k}{q}{-8}%
  \bleft{c}{}\babove{p}{}\bbelow{q}{}\bright{k}{}}}

\vspace{12pt}
\ealpara{Wypełnij luki.}{\instruction{Fill in the gaps.}}

\begin{gapq}{0.5}{10}{4xy+2z-5}
\ealpara{To \ealkeytr{wyrażenie} ma \gap\ \ealkeytr{wyrazy}.}{This \ealkey{expression} has \gap\ \ealkey{terms}.}

\ealpara{Jego \ealkeytr{zmienne} to \lgap, \lgap\ i \lgap.}{Its \ealkey{variables} are \lgap, \lgap\ and \lgap.}

\ealpara{\ealkeytr{Współczynnik} przy $xy$ wynosi \gap.}{The \ealkey{coefficient} of $xy$ is \gap.}

\ealpara{\ealkeytr{Stała} wynosi \gap.}{The \ealkey{constant} is \gap.}
\end{gapq}\hfill
\begin{gapq}{0.47}{11}{6m-n+2p+8}
\ealpara{To \ealkeytr{wyrażenie} ma \gap\ \ealkeytr{wyrazy}.}{This \ealkey{expression} has \gap\ \ealkey{terms}.}

\ealpara{\ealkeytr{Współczynnik} przy $m$ wynosi \gap.}{The \ealkey{coefficient} of $m$ is \gap.}

\ealpara{\ealkeytr{Współczynnik} przy $n$ wynosi \gap.}{The \ealkey{coefficient} of $n$ is \gap.}

\ealpara{$8$ nazywa się \wgap.}{The $8$ is called the \wgap.}
\end{gapq}

\vspace{12pt}
\begin{gapq}{0.5}{12}{-3ab+7}
\ealpara{To \ealkeytr{wyrażenie} ma \gap\ \ealkeytr{wyrazy}.}{This \ealkey{expression} has \gap\ \ealkey{terms}.}

\ealpara{Jego \ealkeytr{zmienne} to \lgiven{$a$} i \lgap.}{Its \ealkey{variables} are \lgiven{$a$} and \lgap.}

\ealpara{\ealkeytr{Współczynnik} przy $ab$ wynosi \gap.}{The \ealkey{coefficient} of $ab$ is \gap.}

\ealpara{\ealkeytr{Stała} wynosi \gap.}{The \ealkey{constant} is \gap.}
\end{gapq}\hfill
\begin{gapq}{0.47}{13}{2a+b-4c+d-9}
\ealpara{To \ealkeytr{wyrażenie} ma \given{$5$} \ealkeytr{wyrazów}.}{This \ealkey{expression} has \given{$5$} \ealkey{terms}.}

\ealpara{\ealkeytr{Współczynnik} przy $b$ wynosi \gap.}{The \ealkey{coefficient} of $b$ is \gap.}

\ealpara{\ealkeytr{Współczynnik} przy $c$ wynosi \gap.}{The \ealkey{coefficient} of $c$ is \gap.}

\ealpara{\ealkeytr{Stała} wynosi \gap.}{The \ealkey{constant} is \gap.}
\end{gapq}

\vspace{12pt}
\ealpara{\ealkeytr{Wyrażenie} ma \ealkeytr{zmienne} $p$ i $q$. \ealkeytr{Współczynnik} przy $p$ wynosi $3$, \ealkeytr{współczynnik} przy $q$ wynosi $-5$, a \ealkeytr{stała} wynosi $4$. Zapisz \ealkeytr{wyrażenie}.}{\qn{14} An \ealkey{expression} has the \ealkey{variables} $p$ and $q$. The \ealkey{coefficient} of $p$ is $3$, the \ealkey{coefficient} of $q$ is $-5$ and the \ealkey{constant} is $4$. Write down the \ealkey{expression}.}
\answerline{3.4cm}

\vspace{12pt}
\ealpara{Zapisz \ealkeytr{wyrażenie}, które ma $4$ \ealkeytr{wyrazy}, $2$ \ealkeytr{zmienne} i \ealkeytr{stałą} równą $-1$.}{\qn{15} Write an \ealkey{expression} that has $4$ \ealkey{terms}, $2$ \ealkey{variables} and a \ealkey{constant} of $-1$.}
\answerline{5.5cm}

\newpage

\ealpara{Wpisz nazwę każdej części w jej pole.}{\instruction{Write the name of each part in its box.}}

\labelq{16}{\labelled{k}{%
  \pfirst{k}{8}\pterm{c}{k}{-3}\pnext{x}{c}{x}%
  \bleft{k}{}\babove{c}{}\bright{x}{}\brbelow{k}{x}{}}}

\vspace{14pt}
\labelq{17}{\labelled{c}{%
  \pfirst{c}{2}\pnext{a}{c}{a}\pterm{k}{a}{-7}\pterm{d}{k}{+0.5}\pnext{b}{d}{b}%
  \bleft{c}{}\babove{k}{}\bright{b}{}\brbelow{d}{b}{}}}

\vspace{14pt}
\labelq{18}{\labelled{c}{%
  \pfirst{c}{-4}\pnext{m}{c}{m}\pnext{n}{m}{n}\pterm{k}{n}{+1.5}%
  \bleft{c}{}\babove{m}{}\brbelow{c}{n}{}\bright{k}{}}}

\vspace{12pt}
\ealpara{Wypełnij luki.}{\instruction{Fill in the gaps.}}

\begin{gapq}{0.5}{19}{0.5a-b+2.5+3c}
\ealpara{To \ealkeytr{wyrażenie} ma \gap\ \ealkeytr{wyrazy}.}{This \ealkey{expression} has \gap\ \ealkey{terms}.}

\ealpara{\ealkeytr{Współczynnik} przy $a$ wynosi \gap.}{The \ealkey{coefficient} of $a$ is \gap.}

\ealpara{\ealkeytr{Współczynnik} przy $b$ wynosi \gap.}{The \ealkey{coefficient} of $b$ is \gap.}

\ealpara{\ealkeytr{Stała} wynosi \gap.}{The \ealkey{constant} is \gap.}
\end{gapq}\hfill
\begin{gapq}{0.47}{20}{9-0.5k^2}
\ealpara{To \ealkeytr{wyrażenie} ma \gap\ \ealkeytr{wyrazy}.}{This \ealkey{expression} has \gap\ \ealkey{terms}.}

\ealpara{Jego \ealkeytr{zmienna} to \lgap.}{Its \ealkey{variable} is \lgap.}

\ealpara{\ealkeytr{Stała} wynosi \gap.}{The \ealkey{constant} is \gap.}

\ealpara{\ealkeytr{Współczynnik} przy $k^2$ wynosi \gap.}{The \ealkey{coefficient} of $k^2$ is \gap.}
\end{gapq}

\vspace{12pt}
\begin{gapq}{0.5}{21}{2xy-5x+0.3y}
\ealpara{To \ealkeytr{wyrażenie} ma \gap\ \ealkeytr{wyrazy}.}{This \ealkey{expression} has \gap\ \ealkey{terms}.}

\ealpara{\ealkeytr{Współczynnik} przy $xy$ wynosi \gap.}{The \ealkey{coefficient} of $xy$ is \gap.}

\ealpara{$-5$ to \ealkeytr{współczynnik} przy \gap.}{$-5$ is the \ealkey{coefficient} of \gap.}

\ealpara{To \ealkeytr{wyrażenie} nie ma \wgap.}{This \ealkey{expression} has no \wgap.}
\end{gapq}\hfill
\begin{gapq}{0.47}{22}{x^2+x-y+xy-3}
\ealpara{To \ealkeytr{wyrażenie} ma \given{$5$} \ealkeytr{wyrazów}.}{This \ealkey{expression} has \given{$5$} \ealkey{terms}.}

\ealpara{\ealkeytr{Współczynnik} przy $x^2$ wynosi \gap.}{The \ealkey{coefficient} of $x^2$ is \gap.}

\ealpara{\ealkeytr{Współczynnik} przy $xy$ wynosi \gap.}{The \ealkey{coefficient} of $xy$ is \gap.}

\ealpara{\ealkeytr{Stała} wynosi \gap.}{The \ealkey{constant} is \gap.}
\end{gapq}

\vspace{12pt}
\ealpara{\ealkeytr{Wyrażenie} ma \ealkeytr{zmienne} $x$, $y$ i $z$. \ealkeytr{Współczynnik} przy $x$ wynosi $-1$, \ealkeytr{współczynnik} przy $y$ wynosi $0.5$, \ealkeytr{współczynnik} przy $z$ wynosi $2$, a \ealkeytr{stała} wynosi $-3$. Zapisz \ealkeytr{wyrażenie}.}{\qn{23} An \ealkey{expression} has the \ealkey{variables} $x$, $y$ and $z$. The \ealkey{coefficient} of $x$ is $-1$, the \ealkey{coefficient} of $y$ is $0.5$, the \ealkey{coefficient} of $z$ is $2$ and the \ealkey{constant} is $-3$. Write down the \ealkey{expression}.}
\answerline{3.4cm}

\vspace{12pt}
\ealpara{Zapisz \ealkeytr{wyrażenie}, które ma $5$ \ealkeytr{wyrazów}, $3$ \ealkeytr{zmienne} i \ealkeytr{stałą} równą $0.5$, gdzie każdy \ealkeytr{współczynnik} jest \ealkeytr{ujemny}.}{\qn{24} Write an \ealkey{expression} that has $5$ \ealkey{terms}, $3$ \ealkey{variables} and a \ealkey{constant} of $0.5$, where every \ealkey{coefficient} is \ealkey{negative}.}
\answerline{5.5cm}

\end{document}